\documentclass[conference]{IEEEtran}
\IEEEoverridecommandlockouts
% AutoInspect-X — Capstone literature review.
% Author block is a placeholder to be replaced before any submission.

\usepackage{cite}
\usepackage{amsmath,amssymb,amsfonts}
\usepackage{textcomp}
\usepackage{xcolor}
\usepackage{booktabs}
\def\BibTeX{{\rm B\kern-.05em{\sc i\kern-.025em b}\kern-.08em
    T\kern-.1667em\lower.7ex\hbox{E}\kern-.125emX}}

\begin{document}

\title{Vehicle Exterior Damage Segmentation: A Review of Convolutional,
Transformer, and Hybrid Architectures for Automotive Damage
Inspection\\
{\footnotesize AutoInspect-X Capstone Literature Review}
}

\author{\IEEEauthorblockN{Student Name (placeholder)}
\IEEEauthorblockA{\textit{Department of Computer Science} \\
\textit{Institution} \\
City, Country \\
email@example.com}
}

\maketitle

\begin{abstract}
Automated visual inspection of vehicle exterior damage is a long-standing
application of computer vision for insurance claims, used-vehicle inspection,
and driving-assistance systems. The task has two dominant formulations ---
damage detection/classification and damage \emph{segmentation}. A 2025
systematic review confirms that detection and classification have a large,
saturated applied literature, whereas pixel-level segmentation of fine,
spread-out damage remains an active front. This review synthesizes verified
literature across four threads relevant to a bounded, single-photo damage
\textbf{segmentation} study: (i) public damage-segmentation datasets (CarDD,
VehiDE, CrashCar101 and related resources); (ii) convolutional segmentation
architectures (FCN, U-Net, ResNet, DeepLab-style atrous designs, HRNet,
Mask~R-CNN) and efficient networks; (iii) transformer and CNN--transformer
hybrid designs (ViT, SETR, SegFormer, Swin, TransUNet) and the role of
attention for global context; and (iv) uncertainty/confidence estimation
(aleatoric vs.\ epistemic taxonomy, MC dropout, deep ensembles) as the basis of
an honesty contract between model predictions and claims. We close with the
research gaps this evidence supports and clearly-separate \emph{literature
findings} from the \emph{preliminary experimental observations} of the
AutoInspect-X capstone, whose two 15-epoch pilots are explicitly not a
conclusion. All cited works were verified against publisher, arXiv, or project
pages in September 2026 unless flagged; no claim is strengthened beyond what
the cited sources support.
\end{abstract}

\begin{IEEEkeywords}
vehicle damage, semantic segmentation, CNN, vision transformer, CNN--transformer
hybrid, encoder--decoder, attention, mean IoU, uncertainty, confidence
calibration.
\end{IEEEkeywords}

\section{Introduction}
\IEEEPARstart{T}{he} automatic assessment of vehicle exterior damage has been a
practical computer vision problem for decades: an insurer photographing a claim,
a used-car inspector, or a driver documenting an incident all need the same
basic question answered --- \emph{where is the damage and how bad is it?} The
commercial value is large: automated triage can route the simplest claims to
low-cost settlement and free human assessors for complex cases. The most recent
systematic review of the field screened 55 selected papers and grouped methods
by task (classification, detection, segmentation) and application context
(insurance claim, traffic, post-accident) \cite{hasan2025}. The picture it
paints is consistent: damage \textbf{detection} (bounding boxes) and damage
\textbf{classification} (scratch vs.\ dent vs.\ crack) are mature and saturated;
damage \textbf{segmentation} (assigning a class label to each damaged pixel) is
where the field is actively moving, because insurance and inspection applications
need shape, extent, and location, not just a label. Commercial systems are
closed-source and do not publish evaluation curves \cite{hasan2025}.

Yet the honest research value is bounded by data. Publicly available datasets
with pixel-level damage annotations exist (CarDD \cite{cardd2023}, VehiDE
\cite{vehide2023}, CrashCar101 \cite{crashcar101}), but no public, licence-clear
dataset carries observed repair \emph{costs}; cost prediction therefore depends
on private or synthetic sources and is not reproducible ground truth.

Three threads in the wider segmentation literature are directly relevant.
First, \emph{convolutional} encoder--decoder designs (U-Net \cite{unet2015})
remain the workhorse for small-data, fine-structure segmentation and are the
ancestor of most automotive damage segmentation baselines. Second,
\emph{transformer} and \emph{CNN--transformer hybrid} designs
\cite{vit2021,setr2021,segformer2021,transunet2021,swin2021} add global context
that plain convolutions struggle to obtain, which matters when a dent on a door
and a displaced panel on an A-pillar are physically far apart yet causally
linked. Third, \emph{uncertainty estimation} \cite{kendall2017,mcdropout2016,
deepensembles2017} supplies the vocabulary for a system that labels its own
confidence honestly instead of emitting raw softmax values as if they were
truth.

This review serves a specific capstone --- \textbf{AutoInspect-X}, a photo-first
car damage \emph{segmentation} demo with an explicit research question: does a
CNN+transformer hybrid achieve higher segmentation quality than a plain U-Net
baseline on CarDD under identical splits, seed, and schedule, and how honestly
can the models label their own confidence? We restrict the review to verified,
reproducible sources; claims are separated into what a source states, what is
inferred from the evidence, and what is a preliminary observation of the
capstone's own (single-seed, 15-epoch) pilots.

The review is organized as follows. Section~\ref{sec:scope} states the selection
and integrity method. Section~\ref{sec:metrics} defines the segmentation task and
the metrics used to compare methods. Section~\ref{sec:data} reviews the verified
damage datasets. Sections~\ref{sec:cnn}--\ref{sec:hybrid} review the three
architectural families (convolutional, transformer, hybrid). Section~\ref{sec:uncertainty}
covers uncertainty and confidence. Section~\ref{sec:gaps} synthesizes the gaps,
and Section~\ref{sec:pilots} separates literature findings from the capstone's
own preliminary observations.

\section{Scope and Selection Method}
\label{sec:scope}
We follow the selection criterion adopted in the AutoInspect-X research report:
a work is included if it (i) defines the dataset, baseline, or training rule of
this project; (ii) is a directly comparable architecture for the research
question (transformer and hybrid segmentation); or (iii) supplies the
uncertainty vocabulary that the \emph{honesty contract} relies on. This is a
targeted review for a bounded capstone, not an exhaustive survey; readers
interested in a systematic landscape are directed to the systematic review of
Hasan et al.~\cite{hasan2025}.

Two integrity rules are enforced throughout. First, every citation was checked
against publisher, arXiv, or project pages during the September 2026 revision of
the project research report; where a previously-listed dataset could not be
verified it was excluded (this applies to the unverified ``CDD (2025)'' and
``Insurance-Damage-v2'' names, which are not used as evidence here). Second, we
never treat cross-paper numbers as comparable benchmarks: the studies differ in
dataset version, split policy, pre-processing, and evaluation code, so their
reported numbers are not directly comparable \cite{cardd2023,hasan2025}.

\section{Segmentation Task and Evaluation Metrics}
\label{sec:metrics}
\subsection{Task formulations}
Automotive damage analysis appears in four formulations that recur across the
literature \cite{hasan2025}:
\begin{itemize}
  \item \textbf{Classification} --- one label per image or per region
  (scratch, dent, crack, glass shatter, lamp broken, tire flat). Saturated.
  \item \textbf{Detection} --- bounding boxes per damage instance. Mature.
  \item \textbf{Semantic / instance segmentation} --- per-pixel class labels
  (semantic) or per-instance masks (instance). The active front
  \cite{cardd2023,vehide2023,crashcar101}.
  \item \textbf{Salient object detection} --- separating ``damaged region'' from
  background without fine classes; supported by CarDD and VehiDE
  \cite{cardd2023,vehide2023}.
\end{itemize}
AutoInspect-X targets \emph{semantic segmentation} of damage classes from a
single photo, which is the smallest reproducible formulation that still answers
``where is the damage and how extensive''.

\subsection{Metrics}
The standard segmentation metrics descend from the PASCAL VOC evaluation
protocol \cite{voc2010}. Mean intersection-over-union (mIoU) over classes is the
primary metric in most damage-segmentation studies \cite{cardd2023,vehide2025}.
For $C$ classes, with ground-truth pixels $G_c$ and predicted pixels $P_c$ for
class $c$:
\begin{equation}
\mathrm{IoU}_c = \frac{|G_c \cap P_c|}{|G_c \cup P_c|},
\qquad
\mathrm{mIoU} = \frac{1}{C}\sum_{c=1}^{C}\mathrm{IoU}_c .
\label{eq:iou}
\end{equation}
The Dice coefficient, closely related to IoU,
\begin{equation}
\mathrm{Dice}_c = \frac{2|G_c \cap P_c|}{|G_c| + |P_c|} ,
\label{eq:dice}
\end{equation}
rewards overlap slightly differently (it equals the harmonic mean of
precision and recall) and is often reported alongside IoU
\cite{voc2010,cardd2023}. Pixel accuracy --- the fraction of all pixels correctly
classified --- is simple to compute but is a weak summary when the background
dominates, which is always true in damage photos.

Four practical warnings about segmentation metrics recur in the literature
\cite{voc2010,hasan2025} and are adopted by AutoInspect-X:
\begin{enumerate}
  \item \textbf{Background dominance.} If the background class is included in
  the mean, a large, correctly-predicted background can inflate mIoU; the
  AutoInspect-X protocol therefore reports \emph{foreground-only} means, and
  absent classes score zero.
  \item \textbf{Class imbalance.} Most damage classes are rare at the pixel
  level (a crack is a few thousand pixels in a $512\times512$ image). Mean
  metrics over classes therefore weigh rare classes equally with common ones,
  which is desirable but must be stated.
  \item \textbf{Small-damage slice.} The binding constraint for trust in real
  claim photos is not the overall mean but the IoU restricted to the
  \emph{smallest} ground-truth instances. The AutoInspect-X protocol defines a
  small-damage slice (damage instances at or below the 25th percentile of the
  training-set instance area) and reports slice IoU/Dice explicitly.
  \item \textbf{Split discipline.} Cross-paper numbers are not comparable because
  papers differ in dataset version, split policy, pre-processing, and evaluation
  code \cite{hasan2025}; within one study, the test split must be touched only
  once after model selection.
\end{enumerate}
\subsection{Task formulations in AutoInspect-X}
A single-damage photo may contain several instances of several classes. The
AutoInspect-X protocol uses the CarDD instance masks as \emph{real ground
truth}, collapses them to a per-pixel semantic label (background or one of the
six classes) at $512\times512$, and evaluates foreground mIoU \eqref{eq:iou}, Dice
\eqref{eq:dice}, and pixel accuracy on the official train/val/test splits.
\section{Vehicle Damage Datasets}
\label{sec:data}
Table~\ref{tab:data} summarizes the verified public datasets used across damage
analysis work that informed this review.

\begin{table}[!t]
\setlength{\tabcolsep}{4.5pt}
\renewcommand{\arraystretch}{1.15}
\caption{Verified public vehicle-damage datasets relevant to segmentation.}
\label{tab:data}
\centering
\begin{tabular}{@{}p{1.9cm}p{1.1cm}p{1.6cm}p{2.6cm}@{}}
\toprule
\textbf{Dataset} & \textbf{Images} & \textbf{Annotation} & \textbf{Verified?}\\
\midrule
CarDD \cite{cardd2023} & 4,000 & Polygon instance masks; 6 classes (dent, scratch, crack, glass shatter, lamp broken, tire flat) & Yes\\
VehiDE \cite{vehide2023} & 13,945 & Detection, segmentation, salient object detection; 8 categories & Yes\\
\texttt{CrashCar101} \cite{crashcar101} & 101,050 & Synthetic mask; part + damage, severity & Yes\\
CDA-Net data \cite{cdanet2023} & 11,380 & Damage annotations, Indian market & Yes\\
Three-quarter-view \cite{threequarter2024} & n/a & Severity classification, \emph{Heliyon} & Yes\\
\bottomrule
\end{tabular}
\end{table}

\subsection{CarDD --- working dataset of this project}
CarDD (Wang, Li, Wu, 2023) is the first publicly available large-scale dataset
for vision-based car damage detection and segmentation: 4,000 high-resolution
images with over 9,000 annotated instance masks across the six damage classes
above \cite{cardd2023}. It supports object detection, instance segmentation, and
salient object detection, and its primary reference model in the original paper
is Mask~R-CNN \cite{cardd2023,maskrcnn2017}. For AutoInspect-X, CarDD is used
with its official splits preserved (train 2,816 images / 6,211 instances, val
810 / 1,744, test 374 / 785 --- measured from a local audit). CarDD contains
damage masks only: no vehicle-part masks, no vehicle metadata, and no repair
cost or repair-action annotations. This is why cost prediction is out of scope
for the project: there is no public cost ground truth to learn from.

\subsection{VehiDE}
VehiDE \cite{vehide2023} is a publicly released large-scale resource for
segmenting and detecting automotive damage: 13,945 high-resolution photos with
more than 32,000 annotated damage instances across eight damage categories. A
follow-up study benchmarks segmentation methods on it
\cite{vehide2025}. An important verification note from the report: an earlier
revision of the project documentation listed ``VeHIDE'' as ``12,000 images, 61
parts, 26 damage types''; that figure is \textbf{not} supported by the verified
VehiDE numbers and instead matches the dataset described by the ALBERT model
paper \cite{albert2025}. The 12,000/61/26 figure is treated as unverified and is
not used as evidence.

\subsection{CrashCar101 and other resources}
CrashCar101 \cite{crashcar101} is a procedurally generated dataset of 101,050
car-crash images with \emph{synthetic} part and damage segmentation masks and
structural severity levels. Because its imagery is synthetic, it complements
real datasets as an augmentation or pre-training source, not as real-world
validation evidence. CDA-Net \cite{cdanet2023} introduces an 11,380-image damage
dataset from the Indian market together with a YOLOv5-based severity pipeline
(when published, YOLOv5 itself preceded the YOLOv8 software release discussed in
Section~\ref{sec:detection}). The three-quarter-view dataset
\cite{threequarter2024} is a severity-classification resource for whole-car-side
photos. A multi-view detection model (MVA-CDD, Peng et al., 2025)
\cite{mvacdd2025} confirms active research interest in multi-view fusion, which
is beyond this project's single-photo scope.

\subsection{Dataset-family synthesis}
Four dataset families recur: (i) \emph{real, segmentation-annotated} (CarDD,
VehiDE) --- verifiable and licence-clear, the only suitable training source for
this capstone; (ii) \emph{synthetic/procedural with perfect masks}
(CrashCar101) --- a training complement, not validation evidence; (iii)
\emph{severity classification} (three-quarter-view; small severity sets); and
(iv) \emph{closed commercial} (unreleased insurer data, not verifiable).

\subsection{Why CarDD is the working dataset}
Two properties make CarDD \cite{cardd2023} the defensible working dataset for a
bounded capstone with a 4~GB GPU. First, it is \emph{verifiable}: the dataset and
its splits are public, and its annotation schema (polygon instance masks over
six damage classes) can be audited locally; the AutoInspect-X local audit
confirms train 2,816/~6,211 instances, val 810/~1,744, test 374/~785. These
instance counts document a further subtlety: the dataset is \emph{class
imbalanced} both at image level (dent and scratch predominate) and pixel level,
so the metrics of \eqref{eq:iou}--\eqref{eq:dice} are reported per class as well
as aggregated. Second, it is \emph{representative of the regime} the demo
targets: uncontrolled photographs with thin, spread-out, partly occluded damage.
VehiDE is a natural future extension for cross-dataset generalization checks
\cite{vehide2023,vehide2025}, and CrashCar101 \cite{crashcar101} could support
pre-training or augmentation --- both are explicitly deferred rather than adopted,
because training on additional data without a documented, licence-reviewed
protocol would violate the reproducibility rules of the project. Synthetic
masks, in particular, are a complement for training but never validation
evidence \cite{crashcar101}.

\subsection{Small-damage framing across datasets}
All three primary datasets are dominated by small damage instances in realistic
photos --- a localized dent or a scratch occupies a tiny fraction of the frame.
The AutoInspect-X protocol therefore treats the small-damage slice as the
binding constraint (Section~\ref{sec:metrics}, warning 3) rather than as an
afterthought, and reports slice metrics separately from aggregate metrics. This
is an explicitly-stated project convention; the literature itself does not
agree on a standard small-damage definition \cite{hasan2025}.

\section{Convolutional Architectures for Segmentation}
\label{sec:cnn}
\subsection{Fully convolutional networks and ResNets}
Semantic segmentation in deep learning begins with the fully convolutional
network (FCN), which replaces the fully-connected layers of a classifier with
transposed-convolution upsampling so the network is trainable end-to-end on
dense labels \cite{fcn2015}. Backbone capacity is typically provided by residual
networks (ResNet), whose identity skip connections enable very deep networks
\cite{resnet2016}. Both are the canonical building blocks of damage-segmentation
pipelines: a classification backbone that has been turned into a dense-label
decoder.

\subsection{FCN and ResNet in the encoder--decoder pipeline}
\begin{figure}[!t]
\begin{center}
\begin{tabular}{@{}c@{\ }c@{\ }c@{\ }c@{\ }c@{}}
RGB input & $\rightarrow$ & Encoder & $\rightarrow$ & mask \\
$512\times512$ & & 4 stages & & $512\times512$ \\
 & & (channel \\ & & depth grows) & & 7 classes \\
\end{tabular}
\end{center}
\caption{Broad dataflow of the AutoInspect-X segmentation models: a
fully-convolutional encoder--decoder with a bottleneck, evaluated by
\eqref{eq:iou}--\eqref{eq:dice}.}
\label{fig:fcn}
\end{figure}

\subsection{U-Net}
The U-Net \cite{unet2015} is the canonical small-data segmentation architecture:
an encoder that progressively downsamples while growing channel depth, a
bottleneck, and a decoder that upsamples back to the input resolution, with
\emph{skip connections} carrying high-resolution detail from encoder to decoder.
In its original formulation each stage is a pair of $3\times3$ convolutions
(double-conv blocks), downsampling is max-pool, and upsampling is a transposed
convolution; the final layer is a $1\times1$ convolution mapping features to logits
per pixel. It was designed for biomedical data with limited labels and has become
the de-facto baseline for damage segmentation because damage photos share that
regime: few images, small and thin structures of interest. The AutoInspect-X
baseline (ResNet34UNet, and its earlier CarddUNet predecessor) is a direct
descendant of this design: the encoder stacks ResNet-derived blocks, and the
decoder mirrors them with skip connections that concatenate encoder features to
each upsampling stage.

\subsection{Atrous designs and high-resolution representations}
DeepLab-style designs generalize convolutions with \emph{atrous} (dilated)
rates to control the resolution of feature maps and capture multi-scale context
without losing spatial detail \cite{deeplabv3}. HRNet instead maintains
high-resolution representations throughout the network by fusing parallel
multi-resolution branches, explicitly preserving the spatial detail needed for
thin, spread-out damage \cite{hrnet2020}. For efficiency (relevant to the 4~GB
GPU constraint of this capstone), MobileNetV2-style inverted-residual bottlenecks
reduce FLOPs and parameters while retaining representational capacity
\cite{mobilenetv2}.

\subsection{Instance segmentation: Mask R-CNN}
Mask~R-CNN extends the Faster-RCNN detection paradigm with a parallel mask
branch, producing an instance mask per detected object \cite{maskrcnn2017}. It is
the reference model of the CarDD paper \cite{cardd2023} and the standard
instance-level baseline in damage work. Instance masks separate individual dents
or scratches; semantic labels say nothing about which instance a pixel belongs
to. For damage severity and counting, instance-level information is useful, but
it requires bounding-box supervision and is heavier than pure semantic
segmentation.

\subsection{Detection frameworks used in applied damage work}
\label{sec:detection}
Applied damage pipelines frequently use YOLO-family detectors. Ultralytics
YOLOv8 \cite{yolov8} is open-source detection and segmentation software released
in January 2023; it has \emph{no peer-reviewed paper} and must be cited as a
software release. YOLO-style dense detectors are typically trained with a focal
loss that down-weights easy negatives, addressing the extreme background
dominance of dense detection \cite{lin2017focal}. The focal loss
\begin{equation}
\mathrm{FL}(p_t) = -\alpha (1-p_t)^{\gamma}\log(p_t)
\label{eq:focal}
\end{equation}
modulates the standard cross-entropy term by a factor $(1-p_t)^{\gamma}$ that
shrinks the gradient contribution of well-classified examples ($p_t \to 1$) and
focusses training on hard, rare foreground pixels; the balance weight $\alpha$
handles foreground/background class imbalance \cite{lin2017focal}. In
AutoInspect-X, a raw-PyTorch U-Net/hybrid is preferred over a YOLO pipeline to
keep the loss, metrics, and class decoding under control on identical splits and
seeds; a focal-loss variant is a listed alternative that has not been
adopted, and its use would be recorded in an experiment decision rather than
assumed.

\subsection{Small-damage-specific considerations}
Two mechanisms from the reviewed CNN literature are especially relevant to
thin, spread-out damage: \emph{atrous} (dilated) convolution and
\emph{high-resolution fusion}. Atrous convolution inserts dilation into the
kernel so that a fixed-size filter covers a larger receptive field without
extra parameters, letting multi-scale context be captured at full resolution
\cite{deeplabv3}. HRNet carries a high-resolution branch through the whole
network and fuses parallel streams, explicitly preserving the spatial detail
that thin scratches and hairline cracks occupy \cite{hrnet2020}. Both ideas
address the same failure mode that dominates under-trained damage models: the
model collapses small structures into background. In the AutoInspect-X protocol
this failure is measured directly by the small-damage slice (Section~\ref{sec:metrics},
warning 3).

\subsection{CNNs, transformers, hybrids: position in this review}
\begin{table}[!t]
\setlength{\tabcolsep}{4.5pt}
\renewcommand{\arraystretch}{1.15}
\caption{Representative architectures and their role in a damage-segmentation
study. Architectural properties are stated as in the cited works; no benchmark
comparison is implied.}
\label{tab:arch}
\centering
\begin{tabular}{@{}p{2.1cm}p{1.9cm}p{3.4cm}@{}}
\toprule
\textbf{Family} & \textbf{Representative} & \textbf{Relevant property}\\
\midrule
CNN & FCN \cite{fcn2015}; ResNet \cite{resnet2016} & Dense-label decoding; identity skips\\
Encoder--decoder CNN & U-Net \cite{unet2015} & Skips preserve fine detail; canonical baseline\\
Multi-scale (atrous) & DeepLab \cite{deeplabv3}; HRNet \cite{hrnet2020} & Atrous dilation; high-res fusion\\
Instance & Mask R-CNN \cite{maskrcnn2017} & Per-instance masks; CarDD's reference model\\
Detection & YOLOv8 \cite{yolov8}; focal loss \cite{lin2017focal} & Applied pipelines; imbalance handling\\
Transformer encoder & ViT \cite{vit2021}; SETR \cite{setr2021} & Global context; patch sequences\\
Hierarchical transformer & Swin \cite{swin2021}; SegFormer \cite{segformer2021} & Multi-scale; linearized attention\\
Hybrid & TransUNet \cite{transunet2021} & Bottleneck attention + CNN decoder\\
Damage-specific & Mask R-CNN on CarDD \cite{cardd2023}; ALBERT \cite{albert2025}; C-DiffDet+ \cite{cdiffdet2025} & Domain adaptations\\
\bottomrule
\end{tabular}
\end{table}

\subsection{CNNs, transformers, hybrids: role in this study}
The baseline of AutoInspect-X is a U-Net-style encoder--decoder; the proposed
model is a hybrid that inserts a bottleneck transformer. Both are trained with
softmax cross-entropy over the argmax class target, on identical splits, seed,
and schedule, and evaluated by \eqref{eq:iou}--\eqref{eq:dice}. No cross-model
comparative claim is drawn from other papers in this review.

\section{Attention and Global Context}
\label{sec:attention}
Self-attention underlies the modern global-context family. The transformer
architecture \cite{attention2017} computes pairwise relevance between all
positions of a sequence, enabling direct long-range dependencies. Scaled
dot-product attention computes, for query $\mathbf{Q}$, key $\mathbf{K}$, and
value $\mathbf{V}$ matrices,
\begin{equation}
\mathrm{Attention}(\mathbf{Q},\mathbf{K},\mathbf{V})
= \mathrm{softmax}\!\left(\frac{\mathbf{Q}\mathbf{K}^{\mathsf T}}{\sqrt{d_k}}\right)\mathbf{V}
\label{eq:attention}
\end{equation}
where $d_k$ is the key dimension; the $\sqrt{d_k}$ scaling stabilizes the
softmax as sequences grow. Multi-head attention runs this in parallel across
several subspaces, allowing the model to attend to different relations
simultaneously \cite{attention2017}. For images, non-local networks graft a
self-attention-style block onto conventional CNNs to capture long-range
dependencies \cite{nonlocal2018}. The Vision Transformer (ViT) \cite{vit2021}
applies a standard transformer encoder directly to \emph{image patches},
treating an image as a sequence of patch tokens and relying on large-scale
pre-training for competitive accuracy.

Two properties matter for damage segmentation. First, attention is
\textbf{quadratic in the number of tokens}: at $H\times W$ spatial resolution the
score matrix $\mathbf{Q}\mathbf{K}^{\mathsf T}$ has $(HW)^2$ entries, which is
why applying global attention to a high-resolution feature map is
prohibitively expensive. Second, attention is a \emph{mechanism for context},
not an architecture: it can be a whole encoder (ViT), a localized window (Swin)
\cite{swin2021}, or a bottleneck module inside a CNN. Where to place it is
therefore a compute--accuracy trade-off, detailed next.

\section{Vision Transformers for Segmentation}
\subsection{ViT as an encoder: SETR}
SETR \cite{setr2021} treats semantic segmentation as a sequence-to-sequence
problem: a ViT encoder consumes the whole image as a sequence of patches, and a
decoder upsamples to full resolution. This gives global context but increases
memory and requires careful training; pure-encoder ViTs also lose the
multi-scale feature hierarchy that CNNs provide for free.

\subsection{Hierarchical transformers: SegFormer and Swin}
Hierarchical transformers --- Swin \cite{swin2021} with shifted windows and
SegFormer \cite{segformer2021} with a positional-encoding-free lightweight
decoder --- restore a CNN-like multi-scale pyramid. Swin restricts global
attention to non-overlapping local windows and \emph{shifts} the window grid
between successive blocks so cross-window relations can propagate layer by
layer; this replaces the quadratic all-pairs cost of
\eqref{eq:attention} with cost that is linear in the number of patches.
SegFormer uses overlapping patch merging, length-agnostic position encoding,
and an all-MLP decoder, unifying semantic and panoptic segmentation. SegFormer
was nominated as a candidate vision backbone in earlier AutoInspect-X
documentation. The trade-off of both hierarchies is realized cost: a full
hierarchical transformer encoder at high resolution is heavy for a 4~GB GPU,
which motivates the hybrid designs below.

\subsection{Hybrids as the pragmatic middle ground}
For a bounded budget, pure-encoder and full-hierarchy definitions bookend the
design space: a full ViT encoder maximizes context at high cost and loses
multi-scale features; a plain CNN minimizes cost with local receptive fields
only. CNN--transformer hybrids split the difference --- the CNN keeps the
high-resolution hierarchy and the transformer adds global context at a single,
cheap resolution. This is the family the next section examines.

\section{CNN--Transformer Hybrid Segmentation}
\subsection{TransUNet}
TransUNet \cite{transunet2021} is the direct methodological antecedent of the
AutoInspect-X hybrid. It keeps a CNN encoder for high-resolution features, adds a
transformer as a \emph{self-attention stage that operates only at the
bottleneck} (the most downscaled feature map), and then decodes with the
original skip connections. This is a low-risk, GPU-bounded upgrade: global
attention is applied where its resolution cost is smallest, while the CNN
decoder and skips preserve fine detail. The hybrid design of this project
(CarddHybrid, now superseded by the architecture-spec v3 \emph{HybridSegmentation})
follows exactly this template.

\subsection{Why bottleneck attention for damage}
Practical constraints support the bottleneck location over a full ViT encoder for
this capstone: (a) VRAM budget --- the observed peak is about 2.85~GB for the
baseline at base~64, and the added transformer capacity must stay within the
4~GB budget; (b) the CNN decoder's skip connections remain necessary to recover
small-damage detail; and (c) the transformer needs only the bottleneck's spatial
resolution (e.g., $32\times32$) to model inter-patch attention, whereas a
full-patch ViT at $224\times224$ costs far more for the same end goal
\cite{vit2021}. These are engineering constraints from the project's own
measurements, not general claims about which location is ``best''.

\subsection{Bottleneck attention in the AutoInspect-X models}
In the architecture-spec~v3 models, the transformer is a compact encoder stage
applied to the $32\times32$ bottleneck feature map: sinusoidal positional
encodings are added to the patch tokens, then residual self-attention layers
(the head count and depth are fixed in the locked configuration) let tokens mix
globally, and the result is projected back to a CNN decoder. Because the
bottleneck token count is small ($32\times32 = 1024$ tokens), the quadratic cost of
\eqref{eq:attention} stays well within the 4~GB budget. The decoder then
upsamples with the original skip connections. This is the TransUNet pattern
\cite{transunet2021} adapted to damage segmentation, and it is the design whose
preliminary pilots are reported in Section~\ref{sec:pilots}.

\subsection{Transformer-based automotive segmentation}
Applying transformers directly to automotive damage, ALBERT \cite{albert2025}
proposes a bidirectional-encoder-based instance segmentation over a large
annotated automotive dataset (26 real damage types, 7 fake-damage artifacts, 61
distinct car parts). It is notable for explicitly disambiguating \emph{fake}/
non-damage artifacts, a real-world nuisance in claim photos. The diffusion-plus-
detector framework C-DiffDet+ \cite{cdiffdet2025} combines generative inpainting
for anomaly scoring with a detection stage and a severity classifier, confirming
a 2024--2025 trend of generative models converging with task-specific detectors.

\section{Uncertainty, Confidence, and Calibration}
\label{sec:uncertainty}
\subsection{Types of uncertainty}
Kendall and Gal \cite{kendall2017} introduced the influential distinction
between \emph{aleatoric} (data-dependent, irreducible) and \emph{epistemic}
(model-dependent, reducible) uncertainty in vision. This taxonomy is the
vocabulary of the AutoInspect-X \emph{honesty contract}: a low-confidence flag
must read as ``the model itself is uncertain'' (epistemic), not as evidence
about the vehicle.

\subsection{Practical estimators}
Two practical epistemic estimators are established: MC dropout \cite{mcdropout2016},
which keeps dropout active at inference and averages stochastic forward passes
as an approximate Bayesian marginalization; and deep ensembles
\cite{deepensembles2017}, which train independently-seeded networks and use the
variance across members as a predictive uncertainty estimate. Both carry a
compute price --- $T$ stochastic forward passes or $M$ networks respectively ---
which is significant on a 4~GB GPU and at demo-time inference. MC dropout is
comparatively cheap to add to an existing checkpoint (dropout must already exist
in the network); deep ensembles require multiple training runs, which is exactly
the multi-seed protocol already planned for the final AutoInspect-X experiment.
Both remain extensions rather than features of the shipped demo. Independent of
the estimator, the integrity-relevant property is \emph{calibration}: claimed
confidence should match observed accuracy. The systematic review
\cite{hasan2025} notes that commercial damage tools do not publish calibration
curves; a system that reports softmax mean confidence is only honest if it also
flags low confidence loudly.

\subsection{The honesty contract}
AutoInspect-X implements a simpler, deterministic honesty signal --- softmax mean
confidence plus a \texttt{low\_confidence} flag below thresholds --- rather than
MC dropout or ensembles, with the stated reason that the demo-grade baseline must
flag itself honestly and often rather than emit confident-looking numbers. The
flag is deliberately conservative: thresholds (\texttt{min\_mean\_confidence},
\texttt{min\_damage\_fraction}) are applied to the deterministic forward pass,
and the engine labels every mask as MODEL PREDICTION even when confidence is
high. MC dropout and deep ensembles remain listed extensions and would be adopted
only via a documented decision. The proposed operational definition of the
research question on confidence honesty is the \emph{separation} between the
confidence distribution over images where the predicted mask agrees with CarDD
ground truth and the distribution where it does not.

\subsection{Calibration as the integrity property}
Regardless of the uncertainty estimator, the integrity-relevant property is
\emph{calibration}: the model is calibrated when its reported confidence matches
its observed accuracy, e.g., images predicted at 90\% mean confidence are correct
about 90\% of the time. Softmax probabilities are well-known to be
over-confident for out-of-distribution inputs, so reporting raw softmax confidence
without a calibration check can mislead a user into treating a guess as a fact.
None of the commercial damage tools surveyed in the systematic review
\cite{hasan2025} publish calibration curves; the honesty contract therefore treats
calibration as a measured, reported property rather than an assumed one.

\section{Research Gaps}
\label{sec:gaps}
Triangulating the reviewed evidence supports three gaps \cite{hasan2025,cardd2023,
vehide2023,albert2025}:
\begin{enumerate}
  \item \textbf{Detection and classification are saturated.} The applied
  literature on damage detection/classification is large; commercial systems are
  closed-source, and calibration is under-reported.
  \item \textbf{Segmentation is the active front.} CarDD, VehiDE, CrashCar101,
  ALBERT, and C-DiffDet+ all indicate that pixel-level and part-level
  segmentation is where the field is moving \cite{cardd2023,vehide2023,
  crashcar101,albert2025,cdiffdet2025}.
  \item \textbf{Confidence honesty is an under-published differentiator.} The
  defensible, data-real opening is segmentation quality plus a verifiable
  confidence honesty contract, because public damage datasets supply masks but
  not financial or vehicle metadata ground truth. Cost prediction, in
  particular, depends on private datasets or synthetic price rules --- neither is
  reproducible truth.
\end{enumerate}
These gaps jointly argue against the two obvious temptations in automotive
damage work. The first is to claim \emph{severity or cost}: neither CarDD nor
any verified public source provides the repair-cost or part-mask ground truth
needed for an honest cost model. The second is to claim \emph{physical extent}:
an uncontrolled photograph has no scale reference, so physical area in
cm$^{2}$ is not defensible; only a normalized image-relative damage-area ratio
is. Both constraints come from the data, not the method, and they define the
boundary of the honesty contract.

\section{Design Implications for the Capstone}
\begin{table}[!t]
\setlength{\tabcolsep}{4.5pt}
\renewcommand{\arraystretch}{1.15}
\caption{Verified literature matrix: representative works, task, and the claim
this review uses. Rows beyond the sourced claim are the review's own
integration, not claims made by the sources.}
\label{tab:matrix}
\centering
\begin{tabular}{@{}p{2.0cm}p{2.0cm}p{3.6cm}@{}}
\toprule
\textbf{Work} & \textbf{Task / method} & \textbf{Claim used here (sourced/verified)}\\
\midrule
CarDD \cite{cardd2023} & detection + instance seg & First public large-scale car-damage seg dataset; 6 classes; 4000 img / 9k+ inst\\
VehiDE \cite{vehide2023,vehide2025} & seg + detection + SOD & Public large-scale damage seg dataset; 13,945 img / 32k+ inst / 8 classes\\
\texttt{CrashCar101} \cite{crashcar101} & part + damage seg & Synthetic masks; complement, not validation\\
CDA-Net \cite{cdanet2023} & detection + severity & Applied YOLOv5 severity pipeline; 11,380-img damage set\\
MVA-CDD \cite{mvacdd2025} & multi-view detection & Active multi-view damage research (out of single-photo scope)\\
Three-quarter view \cite{threequarter2024} & severity classification & Public severity/classification damage resource\\
C-DiffDet+ \cite{cdiffdet2025} & anomaly + detector + severity & 2024--2025 generative + detector trend\\
ALBERT \cite{albert2025} & transformer instance seg & Transformer automotive seg with fake-damage disambiguation\\
U-Net \cite{unet2015} & encoder--decoder seg & Canonical small-data architecture; basis of the baseline\\
DeepLab \cite{deeplabv3} & seg with atrous conv & Multi-scale context at full resolution\\
Mask R-CNN \cite{maskrcnn2017} & instance seg & CarDD paper's reference model\\
HRNet \cite{hrnet2020} & high-res multi-branch & Spatial-detail-preserving backbone\\
SegFormer \cite{segformer2021} & hierarchical transformer seg & Position-free decoder; proposed vision backbone\\
TransUNet \cite{transunet2021} & CNN + transformer hybrid & Bottleneck transformer; the hybrid's antecedent\\
Kendall \& Gal \cite{kendall2017} & uncertainty taxonomy & Aleatoric vs.\ epistemic vocabulary of the honesty contract\\
MC dropout \cite{mcdropout2016} & epistemic UQ & Test-time dropout uncertainty approximation\\
Deep ensembles \cite{deepensembles2017} & epistemic UQ & Ensemble-variance uncertainty\\
Hasan et al. \cite{hasan2025} & systematic review & Confirms saturated detection/classification; rare calibration\\
\bottomrule
\end{tabular}
\end{table}

The reviewed evidence translates into a small number of concrete design choices
for AutoInspect-X, each traceable to a cited source:
\begin{enumerate}
  \item \emph{Pixel-level task, foreground metrics.} Segmentation answers
  ``where and how extensive'' \cite{cardd2023,hasan2025}; foreground-only mIoU
  and Dice avoid background inflation (Section~\ref{sec:metrics}).
  \item \emph{Shared baseline and protocol.} A U-Net-style encoder--decoder
  \cite{unet2015} is the baseline; the hybrid adds a bottleneck transformer
  following TransUNet \cite{transunet2021}. Identical splits, seed, schedule,
  loss, and metric harness for both arms is the only honest way to isolate the
  architecture effect \cite{hasan2025}.
  \item \emph{Attention cost is bounded.} Global attention at the bottleneck
  \eqref{eq:attention} costs $O((HW)^2)$ in tokens, which is affordable at
  $32\times32$ but not at full resolution; the CNN decoder and skips preserve
  the small-damage detail (Sections~\ref{sec:cnn}, \ref{sec:hybrid}).
  \item \emph{Imbalance handling.} Sparse-class damage needs either class
  re-weighting or a focal-style loss \eqref{eq:focal} \cite{lin2017focal};
  whichever is chosen must be recorded as a decision, not defaulted silently.
  \item \emph{Honesty is a measured property.} Confidence is reported with a
  low-confidence flag and the mask is labelled MODEL PREDICTION; calibration and
  separation are planned measurement targets, not assumptions \cite{kendall2017,
  mcdropout2016,deepensembles2017}.
  \item \emph{Reproducibility of records.} Every run records experiment ID,
  dataset version and split, seed, code and model version, hyperparameters,
  metrics, and assumptions --- consistent with the reproducibility expectations
  in the wider damage literature \cite{hasan2025}.
\end{enumerate}
These choices are recommendations of this review; where AutoInspect-X has
departed from a cited practice (for example, preferring a raw-PyTorch model over
a YOLO pipeline, or reporting a deterministic honesty flag rather than MC
dropout), the departure is a documented decision in the companion engineering
documentation, not an accident.

\subsection{Parameter budgets and the 4~GB training constraint}
Efficiency matters in this capstone for a concrete and non-negotiable reason:
the training and demo hardware is a 4~GB GPU. The Reviewed literature supplies
two complementary responses. Architecturally, lightweight designs --- 
MobileNetV2-style inverted residuals \cite{mobilenetv2} --- trade a small amount
of accuracy for a large reduction in FLOPs and parameters, and hierarchical
transformers with windowed attention \cite{swin2021} reduce the quadratic
attention cost. Procedurally, the budget is respected by placing attention only
at the bottleneck and by keeping the input at $512\times512$ for both arms. The
measured engines of the two 15-epoch pilots are within this budget: the baseline
defined and trained to roughly 24.6M parameters and the hybrid to roughly 27.9M
parameters (not counts to be compared against any external checkpoint, but the
two are directly comparable because they share the harness). The observed
training peak for the earlier base-64 baseline was about 2.85~GB at batch~2;
the pilot configuration is run within the same memory ceiling. Reporting these
numbers is a reproducibility requirement \cite{hasan2025}, not an
architecture claim.

\section{Threats to Validation in the Reviewed Literature}
Three validation pitfalls recur across the works surveyed here and are worth
naming explicitly because they are easy to fall into and hard to detect in a
short review. \emph{Split leakage.} Several damage studies do not state whether
near-duplicate images of the same vehicle appear across train and test; the
systematic review flags dataset provenance as a frequent weakness
\cite{hasan2025}. \emph{Benchmark incomparability.} Papers report numbers on
different dataset versions, crops, and evaluation code; treating them as a
leaderboard is unsound. \emph{Overclaiming extent.} Some applied work reports
damage area or severity as if an uncontrolled photograph were a calibrated
measurement; without a scale reference or calibration geometry this is not
defensible. AutoInspect-X treats each of these as a documented protocol
constraint rather than a review finding.

\subsection{Verification discipline for research-facing claims}
Two further disciplines apply because this review feeds a project that operates
under adversarial research-integrity rules. First, the ground-truth policy of
the project distinguishes REAL GROUND TRUTH (measured/verified by a trusted
source, e.g., CarDD masks), WEAK LABEL (noisy/indirect), SYNTHETIC LABEL
(rule/model-generated), DERIVED FEATURE (computed), MODEL PREDICTION (model
output), and ASSUMPTION (team choice); each claim in this document, and each
metric reported by the project, falls into exactly one of these categories. This
is what stops a softmax score in a demo screenshot from being read as verified
damage evidence, and what keeps the 15-epoch pilot numbers in
Section~\ref{sec:pilots} from being read as a research conclusion.
\subsection{Threats to validation in this review}
The review itself has three limitations. First, it is a targeted review for a
bounded capstone, not an exhaustive survey; the systematic review of Hasan
et al.~\cite{hasan2025} should be consulted for coverage. Second, architectural
claims (``bottleneck attention adds global context'') are stated as the
mechanisms described by the cited papers, not as benchmark accomplishments; no
inter-paper comparison is drawn. Third, where the project's own measurements are
reported (parameter counts, VRAM use, pilot metrics), they are single-run
observations from one environment and are explicitly not generalized. These
limitations are intentional and keep the review within what the evidence
supports.

\section{Preliminary Experimental Observations (AutoInspect-X)}
\label{sec:pilots}
\begin{quote}
\textit{The following subsection is \textbf{not} a literature finding; it is the
capstone's own preliminary experimental observation and is included only so that
the reader can relate the review to the companion architecture report. It is
explicitly \emph{not} a conclusion, and no statistical significance is claimed.}
\end{quote}

The project migrated from the demo-era \texttt{CarddHybrid} to the
architecture-spec~v3 research models: a \textbf{baseline} \texttt{ResNet34UNet}
and the proposed \textbf{HybridSegmentation} (CNN encoder $\rightarrow$
transformer bottleneck $\rightarrow$ CNN decoder with skips). Two 15-epoch
\emph{preliminary validation} pilots (seed~0, full official CarDD splits,
identical schedule, softmax cross-entropy over the argmax class target, EMA
weights) were run as an intermediate check:
\begin{itemize}
  \item \texttt{pilot15\_baseline}: foreground mIoU 0.6127, foreground mDice
  0.7440, pixel accuracy 0.8979 (best at epoch 14);
  \item \texttt{pilot15\_hybrid}: foreground mIoU 0.5963, foreground mDice
  0.7320, pixel accuracy 0.8873 (best at epoch 14).
\end{itemize}
Both models were still improving at epoch 14, and neither result establishes
that either architecture is superior: there is a single seed (0), no statistical
test, and no fine-tuned schedule. The planned controlled comparison is a
60-epoch, 3-seed run under the locked protocol, which is required before any
architecture-level claim. The demo default is \texttt{pilot15\_hybrid} per the
project plan, with \texttt{pilot15\_baseline} retained as the controlled
baseline arm. These pilots validate only that both model families train and
evaluate end-to-end on CarDD under the locked schedule; they are recorded in the
project registry together with their hyperparameters and seeds for
reproducibility.

\subsection{Pre-processing and harness}
The AutoInspect-X harness pre-processes CarDD consistently: images are
resized/cropped to $512\times512$ RGB, and the instance masks are collapsed to a
per-pixel class label (background plus the six damage classes), preserving the
official splits. Training applies horizontal flip and mild brightness/contrast
augmentation only; validation and test are unaugmented. The two model arms share
the identical loss (softmax cross-entropy over the argmax class target), the
Adam optimizer, warm-up, cosine decay, EMA weights, and checkpoint rule (best
validation foreground mIoU). Recording this protocol verbatim is what makes any
future claim about one arm versus the other fair \cite{hasan2025}; the pilots
reported in Section~\ref{sec:pilots} both used it.

\subsection{What the small-damage slice shows}
The small-damage slice defined in Section~\ref{sec:metrics} is not only a metric
convention; it is where the practical difficulty concentrates. In the CarDD
training set, the 25th-percentile instance-area threshold is about 3,013.
pixels at $512\times512$ (i.e., roughly 1.1\% of the image), a measured baseline
against which "small" is defined in the project protocol; instances at or below
this threshold are overwhelmingly dents and scratches --- the two most common
damage classes in real claims. Because these classes occur at tens of pixels in
width, a model trained with a plain cross-entropy objective can under-segment
them entirely, collapsing them into background. The pilot runs do not yet claim
to have fixed this: the small-damage slice remains the binding constraint on
trust in both architectures, and it is reported explicitly rather than hidden in
the aggregate. This measurement is project-specific and is included to
illustrate why the review repeatedly stresses fine-structure recovery
mechanisms (U-Net skips, atrous context, high-resolution fusion) rather than
as a benchmark claim.

\section{Conclusion}
Vehicle damage segmentation is a reproducible, data-bounded research problem.
Real, public, mask-annotated datasets exist (CarDD \cite{cardd2023}, VehiDE
\cite{vehide2023}), synthetic resources (CrashCar101 \cite{crashcar101})
complement them for training, and calibration/confidence honesty is
under-published despite being the natural integrity constraint on a photo-first
inspection system. Architecturally, the evidence supports a CNN--transformer
hybrid as a GPU-bounded, contemporary upgrade to a U-Net baseline --- attention
placed at the bottleneck preserves fine detail via the CNN decoder while adding
global context. The AutoInspect-X pilots confirm only that the two model
families train on CarDD under the locked schedule; the architecture comparison
remains to be settled by the planned multi-seed experiment, and the value of the
system hangs on its honesty contract, not on any single metric.

Three statements summarize the position this review leaves the reader with.
\textbf{(1) Task selection determines honesty.} A segmentation task on a
verified, licence-clear dataset (CarDD) is reproducible; severity/cost claims
are not, because the required ground truth does not exist publicly.
\textbf{(2) Architecture is an affordable open question.} The CNN--transformer
hybrid family is contemporary and GPU-bounded; it is answerable in a capstone by
a locked, multi-seed protocol --- no other paper is needed to establish the
comparison. \textbf{(3) Confidence labelling is the under-exploited, measured
contribution.} Reporting a low-confidence flag and measuring
agreement--disagreement separation converts ``the model was confident'' from a
marketing sentence into a testable property \cite{kendall2017}. Everything else
in damage analysis worth doing remains constrained by data that has not yet been
made public.

\section*{Acknowledgment}
This review was prepared as part of the AutoInspect-X capstone. The companion
research report (AutoInspect-X Research Report, September 2026) contains the
full evidence trail, dataset audits, and decision records referenced here.

\section*{References}
\bibliographystyle{IEEEtran}
\bibliography{references}

\end{document}